% !PRINTED 275
\weberpage{275}
To each transformation of degree $n$ of the
form
\weq{2}{
\begin{pmatrix} a, & 0\\ c, & \partial \end{pmatrix}
}
there can then be assigned one and only one transformation of each of the degrees $n'$, $n''$:
\weq{3}{
\begin{pmatrix} a', & 0\\ c', & \partial' \end{pmatrix}, \quad
\begin{pmatrix} a'', & 0\\ c'', & \partial'' \end{pmatrix}
}
and these are determined by the following conditions:
\weq{4}{
\begin{aligned}
a &= a'a'', & \partial &= \partial'\partial'',\\
\partial''c' &\equiv c \pmod{a'}, & \partial'c'' &\equiv c \pmod{a''},
\end{aligned}
}
and conversely, from every pair of transformations of the
form (3) there follows, by (4), one and only one transformation (2). For by
(4) the quantities $\partial'$, $\partial''$ are first determined as the greatest common
divisors of $\partial$ with $n'$ and $n''$, and thereupon $c'$ is determined with
respect to the modulus $a'$, and $c''$ with respect to the modulus $a''$, from the two
last congruences (4).

What now matters above all is to show that the
congruences (4) are preserved when the transformations (2)
and (3), in accordance with \S 69, (8) to (12), by means of the two linear transformations
\[
(\omega, \omega+1), \quad \left(\omega, -\frac{1}{\omega}\right)
\]
are transformed.

Following the formulae mentioned, we set:
\weq{5}{
\begin{aligned}
\begin{pmatrix} a, & 0\\ c, & \partial \end{pmatrix}
\begin{pmatrix} 1, & 0\\ 1, & 1 \end{pmatrix}
&=
\begin{pmatrix} 1, & 0\\ \lambda, & 1 \end{pmatrix}
\begin{pmatrix} a, & 0\\ c_1, & \partial \end{pmatrix}\\
\begin{pmatrix} a', & 0\\ c', & \partial' \end{pmatrix}
\begin{pmatrix} 1, & 0\\ 1, & 1 \end{pmatrix}
&=
\begin{pmatrix} 1, & 0\\ \lambda', & 1 \end{pmatrix}
\begin{pmatrix} a', & 0\\ c_1', & \partial' \end{pmatrix}\\
\begin{pmatrix} a'', & 0\\ c'', & \partial'' \end{pmatrix}
\begin{pmatrix} 1, & 0\\ 1, & 1 \end{pmatrix}
&=
\begin{pmatrix} 1, & 0\\ \lambda'', & 1 \end{pmatrix}
\begin{pmatrix} a'', & 0\\ c_2'', & \partial'' \end{pmatrix}. % [sic] German source prints c_2''; parallelism suggests c_1''.
\end{aligned}
}
\weq{6}{
\begin{aligned}
\begin{pmatrix} a, & 0\\ c, & \partial \end{pmatrix}
\begin{pmatrix} 0, & 1\\ -1, & 0 \end{pmatrix}
&=
\begin{pmatrix} \alpha, & \beta\\ \gamma, & \delta \end{pmatrix}
\begin{pmatrix} a_2, & 0\\ c_2, & \partial_2 \end{pmatrix}\\
\begin{pmatrix} a', & 0\\ c', & \partial' \end{pmatrix}
\begin{pmatrix} 0, & 1\\ -1, & 0 \end{pmatrix}
&=
\begin{pmatrix} \alpha', & \beta'\\ \gamma', & \delta' \end{pmatrix}
\begin{pmatrix} a_2', & 0\\ c_2', & \partial_2' \end{pmatrix}\\
\begin{pmatrix} a'', & 0\\ c'', & \partial'' \end{pmatrix}
\begin{pmatrix} 0, & 1\\ -1, & 0 \end{pmatrix}
&=
\begin{pmatrix} \alpha'', & \beta''\\ \gamma'', & \delta'' \end{pmatrix}
\begin{pmatrix} a_2'', & 0\\ c_2'', & \partial_2'' \end{pmatrix}.
\end{aligned}
}
